\section{Behold your mother!}

\begin{quotex}
Allegory is a rational operation, implying no transition either to a new plane of being or to a new depth of consciousness; it is a figuration, at an identical level of consciousness, of what might very well be known in a different way. The symbol announces a plane of consciousness distinct from that of rational evidence; it is the ``cipher'' of a mystery, the only means of saying something that cannot be apprehended in any other way; a symbol is never ``explained'' once and for all, but must be deciphered over and over again.
\flright{\textsc{Henry Corbin}, \emph{Alone with the Alone}}
\end{quotex}


Allegory is the lowest form of symbolic writing. There is a key -- an isomorphism, as it is called in Category Theory -- that is, a mapping from one concept to another. The purpose of allegory is to simplify the thing symbolized by stripping it of its inessential elements. The allegory is also easier to remember than the thing allegorized.

People may use the word metaphor when they mean allegory. A metaphor is used once to express an idea of vision in more colorful language, for example, ``The moon was a ghostly galleon''. An allegory, on the other hand, is extended.

It is easy to visualize the full moon as a galleon crossing the sky. Allegories like Pilgrim's Progress or Animal Farm are likewise easy to understand. They can be explained.

However, a true symbol defies such analysis and explanation. A mistake is to try to ``explain'' the symbol. Doing so is a way to become popular but not a way to achieve gnosis. A symbol leads into great depths at each level of understanding.

In this sense, Hollywood movies cannot be symbolic. They cannot serve as the means to preserve and transmit traditional symbols, unless accidentally and incidentally. Perhaps some day that will change.

Hence interpreters of symbolism do no more than explain metaphors and allegories. Alleged symbols are just eruptions of the unconscious and in that sense may even elude the director, writer, and actors.

\paragraph{To the Point}


The film mother! is a film that attempts to be deliberately symbolic in the lower sense of the word. Including obvious allusions to biblical stories hardly gives a film any depth, unless perhaps the stories are actually believed. Unfortunately, in our day and age, biblical allusions are not so obvious to everyone. Reviews of the film believe it is a ``metaphor'' or ``allegory'' for something or other. Perhaps to save the environment as the attractive starlet believes.

I shall neither recommend nor dismiss the film. It is a way to pass some time while being confined to the sofa with a sprained ankle, the pain of which exceeds the pain of watching the film to its completion.

The characters have no names since their roles illustrate the allegory. Jennifer Lawrence is the Virgin Mother who is married to a much older man, a popular poet by trade, known by his pronoun ``Him''. We see that his life had become ruined; his idyllic house in the garden had burned down and his creativity had atrophied. As the opposite of Eve, mother restores the lost Paradise, restoring the large old mansion to a pristine condition.

Then stuff happens. Visitors arrive much to Him's joy. The first arrivals are Adam and Eve along with their two sons. Eve feels entitled to live in the mansion, even while initiating its destruction. Cain kills Abel, bring death back into the Paradise recreated by mother. We empathize with her sorrow and sense of helplessness as Him allows event to get out of control.

We learn that the marriage had never been consummated, leaving mother as a Virgin. Desirous of having a child, mother taunts Him. Angered, he takes her by force, although she eventually yields to his will. The next morning she finds herself pregnant and Him becomes prolific. His new book of poetry is an immediate success, attracting ever more visitors.

They are emotionally moved by his beautiful words. That is spot on. How many posts on social media are full of such beautiful quips? We all ``like'' them, as if the act of liking is the same as incarnating the thought. But we see how the lovers of the Poet reveal their true inner selves. They become angry, lustful, crass, violent.

During the commotion mother gives birth to a son. Unable to recognize them for what they are, Him takes the baby to his many admirers. They then kill the baby and eat is corpse as if consuming a holy sacrament.

Understandably distraught, mother sets the mansion ablaze send the interlopers to the hell that they deserve. Mother dies in the conflagration while Him is unscathed. She is idolized and the cycle repeats.


\flrightit{Posted on 2022-07-02 by Cologero}

\begin{center}* * *\end{center}

\begin{footnotesize}
\begin{sffamily}
\texttt{Michael DeJoy on 2022-07-03 at 12:42 said:}

There are two kinds of people: the people of the ``Bright'' and the people of the ``Shade''. The Bright people are the rational, logic, study philosophy and explain allegories to make sense of them. The Shade people understand symbols, give up reading philosophy texts and trying to explain things to Bright people no matter how concerned they appear. Another example of Bright people are those who insist on riding their bicycles as fast as possible thru all intersections thinking how wonderful it all is.

\hfill

\texttt{Juliano on 2022-07-04 at 11:51 said:}

Have you watched Tarkovsky movies? I think there is great symbolism there. Nostalghia and The Sacrifice are my favorites.

\hfill

\texttt{Cologero on 2022-07-04 at 13:44 said:}

Maybe this one about \emph{Stalker} by \textbf{Andrei Tarkovsky}.

\textit{The Zone of Innermost Desires}: \url{https://www.gornahoor.net/?p=16139}


\hfill

\texttt{Juliano on 2022-07-13 at 11:48 said:}

Great article about Stalker, thank you. I still have to watch this one, will do it soon.

\hfill

\texttt{Blade on 2022-08-22 at 19:43 said:}

Have you written a post that scales symbolic terms? Dictionaries are not a great help. I skip over allegory because it's use often appears idiosyncratic to the author. Now I have a clue. Thanks

\hfill

\texttt{Cologero on 2022-08-22 at 22:58 said:}

More or less in this post. That is what most people mean by symbol.

Also \textit{The Elite and the Gospels}\footnote{\url{https://www.gornahoor.net/?p=16286}}.

I can't recall the posts where this is explained: The real symbols are deliberately put into texts by those who know and understand. One of the reasons is so that esoteric teachings are not lost. The people will continue to preserve and publish the texts, even when they don't understand the esoteric meaning of the symbols.

\hfill

\texttt{Blade on 2022-08-22 at 23:17 said:}

Thank you


\hfill

\end{sffamily}
\end{footnotesize}

